\hypertarget{c-extensions-python-c-api-interoperability}{%
\chapter{C Extensions \& Python C-API
Interoperability}\label{c-extensions-python-c-api-interoperability}}

\hypertarget{writing-a-pure-c-extension-module}{%
\subsection{24.1 Writing a Pure C Extension
Module}\label{writing-a-pure-c-extension-module}}

Python allows writing modules directly in C or C++ to access low-level
OS APIs, optimize performance-critical paths, or link with native
hardware drivers.

\hypertarget{minimal-c-extension-boilerplate}{%
\subsubsection{1. Minimal C-Extension
Boilerplate}\label{minimal-c-extension-boilerplate}}

A complete C-extension module requires defining the methods table, the
module descriptor, and the dynamic initialization hook:

\begin{lstlisting}[language=C]
#define PY_SSIZE_T_CLEAN
#include <Python.h>

/* 1. Core C Function Implementation */
static PyObject* method_add(PyObject* self, PyObject* args) {
    long a, b;
    /* Parse Python arguments tuple into native C types */
    if (!PyArg_ParseTuple(args, "ll", &a, &b)) {
        return NULL; /* Raises TypeError automatically */
    }
    long sum = a + b;
    /* Convert native C type back to PyObject integer wrapper */
    return PyLong_FromLong(sum);
}

/* 2. Methods Table Registration */
static PyMethodDef CustomMethods[] = {
    {"add", method_add, METH_VARARGS, "Calculate the sum of two integers."},
    {NULL, NULL, 0, NULL} /* Sentinel marker to end array */
};

/* 3. Module Definition Structure */
static struct PyModuleDef custommodule = {
    PyModuleDef_HEAD_INIT,
    "custom_c_math",            /* Module Name */
    "A custom high-performance C extension module.", /* Docstring */
    -1,                         /* Module state size (-1 = global state) */
    CustomMethods
};

/* 4. Initialization Hook called upon 'import custom_c_math' */
PyMODINIT_FUNC PyInit_custom_c_math(void) {
    return PyModule_Create(&custommodule);
}
\end{lstlisting}

\begin{center}\rule{0.5\linewidth}{0.5pt}\end{center}

\hypertarget{reference-counting-and-ownership-rules-in-c-api}{%
\subsection{24.2 Reference Counting and Ownership Rules in
C-API}\label{reference-counting-and-ownership-rules-in-c-api}}

When interacting with the Python C-API, developers must manually manage
reference counting to avoid memory leaks or segfaults due to premature
deallocation.

\hypertarget{reference-ownership-categories}{%
\subsubsection{1. Reference Ownership
Categories}\label{reference-ownership-categories}}

\begin{itemize}
\tightlist
\item
  \textbf{New References}: The C function receives ownership of the
  reference. It must decrement the reference count when finished, or
  pass ownership back to CPython.

  \begin{itemize}
  \tightlist
  \item
    \emph{Examples}: \passthrough{\lstinline!PyLong\_FromLong()!},
    \passthrough{\lstinline!PyList\_New()!},
    \passthrough{\lstinline!PyObject\_Call()!}.
  \end{itemize}
\item
  \textbf{Borrowed References}: The C function receives a pointer
  without ownership. It does not own the reference and must not
  decrement it unless it explicitly calls
  \passthrough{\lstinline!Py\_INCREF()!} to claim ownership.

  \begin{itemize}
  \tightlist
  \item
    \emph{Examples}: \passthrough{\lstinline!PyTuple\_GetItem()!},
    \passthrough{\lstinline!PyList\_GetItem()!}.
  \end{itemize}
\end{itemize}

\begin{lstlisting}[language=C]
PyObject* list = PyList_New(1);      /* Returns a New Reference */
PyObject* item = PyList_GetItem(list, 0); /* Returns a Borrowed Reference */

Py_INCREF(item);                    /* Converts 'item' to a New Reference */
Py_DECREF(list);                    /* Safely frees list; item remains valid */
Py_DECREF(item);                    /* Frees item */
\end{lstlisting}

\begin{center}\rule{0.5\linewidth}{0.5pt}\end{center}

\hypertarget{interoperability-toolchains-pybind11-vs.-cython}{%
\subsection{24.3 Interoperability Toolchains: Pybind11
vs.~Cython}\label{interoperability-toolchains-pybind11-vs.-cython}}

Writing pure C extensions can be verbose and error-prone. Modern
projects utilize toolchains to automate bindings:

\hypertarget{pybind11}{%
\subsubsection{1. Pybind11}\label{pybind11}}

Pybind11 is a header-only library that exposes C++ types and functions
to Python using advanced template metaprogramming. It generates C-API
wrapper code during compile time, automatically converting types (e.g.,
\passthrough{\lstinline!std::vector!} to Python
\passthrough{\lstinline!list!}).

\hypertarget{cython}{%
\subsubsection{2. Cython}\label{cython}}

Cython is a superset of Python that supports calling C functions and
declaring C types. The Cython compiler translates annotated
\passthrough{\lstinline!.pyx!} files into optimized C code. *
\textbf{Static Type Declarations}: Using \passthrough{\lstinline!cdef!}
keywords avoids Python dynamic lookup overhead by compiling directly to
native C variables and structs. * \textbf{Direct Struct Lookups}: It
accesses C-level structs directly, completely bypassing the dictionary
lookup layer for microsecond-level execution times.

\begin{center}\rule{0.5\linewidth}{0.5pt}\end{center}
